\subsection{Pre-registration and Study Governance}
\label{sec:prereg}

\paragraph{Pre-registration.}
Prior to full-scale packet distribution and rater recruitment, the study design, case inclusion criteria, hypotheses, and primary analysis plan will be pre-registered in a time-stamped public repository (e.g., OSF or equivalent archival platform). The pre-registration document will include:

\begin{enumerate}
    \item The complete list of 25 cases.
    \item The inclusion and exclusion criteria for case selection.
    \item The three-condition generation protocol (Baseline, Protocol, Independent).
    \item The normalization rules and word-budget constraints.
    \item The randomization procedure for document labeling and packet construction.
    \item The primary and secondary outcome definitions.
    \item The planned statistical analysis, including stratified and robustness analyses.
\end{enumerate}

No cases will be added, removed, or replaced after pre-registration except under documented failure of eligibility criteria (e.g., discovery that an independent assessor report does not meet minimum documentation requirements).

\paragraph{Hypotheses.}
The primary pre-registered hypothesis is:

\begin{quote}
Across the full case set, protocol-structured analyses will be selected as closer to independent assessor reports more frequently than baseline analyses.
\end{quote}

A secondary, conditional hypothesis is:

\begin{quote}
The protocol advantage will be larger in governance-heavy, underdetermined cases (R3--R4 dominant) than in mechanically crisp cases (R1--R2 dominant).
\end{quote}

No directional hypothesis is asserted for mechanically crisp cases.

\paragraph{Case Stratification.}
Cases are pre-classified into strata (mechanistic, socio-technical, financial-institutional, regulatory, crisis/disaster). Stratum assignments will be fixed in the pre-registration document and not modified post hoc. Primary aggregate analysis will be performed across all cases; stratum-level analyses will be reported separately.

\paragraph{Generation Freeze.}
All three generated outputs per case (Baseline, Protocol, Independent) will be frozen prior to normalization. No content edits will occur after freezing except formatting normalization consistent with pre-specified rules (Section~\ref{sec:methods}). All raw generations will be archived.

\paragraph{Normalization Protocol Freeze.}
The normalization shell, section headers, word limits, and constraints prohibiting vocabulary injection or reasoning augmentation will be finalized prior to packet assembly. Any deviation from these rules will be documented.

\paragraph{Randomization and Blinding.}
The mapping between document condition (Baseline, Protocol, Independent) and Type labels (X/Y/Z) will be randomized per rater using a fixed seed stored in the administrative archive. Raters will not be informed of the existence of protocol vs baseline conditions.

\paragraph{Primary Analysis Lock.}
The primary statistical test (paired comparison of $(P,I)$ vs $(B,I)$ selection frequency) will be defined prior to inspecting rater responses. Secondary analyses (R3--R4 restricted subset, stratum-level effects) will be reported regardless of outcome.

\paragraph{Full Disclosure Policy.}
The final manuscript will include:
\begin{itemize}
    \item Per-case score tables.
    \item Aggregate and stratum-level results.
    \item Cases where no protocol advantage was observed.
    \item Sensitivity analyses.
\end{itemize}

\paragraph{Exploratory Analyses.}
Any post hoc or exploratory analyses (e.g., textual feature correlation, qualitative pattern analysis) will be clearly labeled as exploratory and separated from confirmatory claims.

\paragraph{Replication Materials.}
Upon publication, all prompts, normalization templates, randomization scripts, anonymized rater responses, and analysis code will be made available to enable independent replication, subject to any confidentiality constraints agreed with raters.
